\documentclass[10pt,a4paper]{amsart}
\usepackage{amsmath,amssymb,amsthm,graphicx}
\usepackage{verbatim}
\newenvironment{reference}{\comment}{\endcomment}
%\newenvironment{reference}{}{}
\newenvironment{slogan}{\comment}{\endcomment}
\newenvironment{history}{\comment}{\endcomment}

% For commutative diagrams we use Xy-pic
\usepackage[all]{xy}

% We use 2cell for 2-commutative diagrams.
\xyoption{2cell}
% Theorem environments.
%
\theoremstyle{plain}
\newtheorem{theorem}[subsection]{Theorem}
\newtheorem{proposition}[subsection]{Proposition}
\newtheorem{lemma}[subsection]{Lemma}

\theoremstyle{definition}
\newtheorem{definition}[subsection]{Definition}
\newtheorem{example}[subsection]{Example}
\newtheorem{exercise}[subsection]{Exercise}
\newtheorem{situation}[subsection]{Situation}

\theoremstyle{remark}
\newtheorem{remark}[subsection]{Remark}
\newtheorem{remarks}[subsection]{Remarks}

\numberwithin{equation}{subsection}

% Macros
%
\def\lim{\mathop{\mathrm{lim}}\nolimits}
\def\colim{\mathop{\mathrm{colim}}\nolimits}
\def\Spec{\mathop{\mathrm{Spec}}}
\def\Hom{\mathop{\mathrm{Hom}}\nolimits}
\def\Ext{\mathop{\mathrm{Ext}}\nolimits}
\def\SheafHom{\mathop{\mathcal{H}\!\mathit{om}}\nolimits}
\def\SheafExt{\mathop{\mathcal{E}\!\mathit{xt}}\nolimits}
\def\Sch{\mathit{Sch}}
\def\Mor{\mathop{\mathrm{Mor}}\nolimits}
\def\Ob{\mathop{\mathrm{Ob}}\nolimits}
\def\Sh{\mathop{\mathit{Sh}}\nolimits}
\def\NL{\mathop{N\!L}\nolimits}
\def\CH{\mathop{\mathrm{CH}}\nolimits}
\def\proetale{{pro\text{-}\acute{e}tale}}
\def\etale{{\acute{e}tale}}
\def\QCoh{\mathit{QCoh}}
\def\Ker{\mathop{\mathrm{Ker}}}
\def\Im{\mathop{\mathrm{Im}}}
\def\Coker{\mathop{\mathrm{Coker}}}
\def\Coim{\mathop{\mathrm{Coim}}}

% Boxtimes
%
\DeclareMathSymbol{\boxtimes}{\mathbin}{AMSa}{"02}

%
% Macros for moduli stacks/spaces
%
\def\QCohstack{\mathcal{QC}\!\mathit{oh}}
\def\Cohstack{\mathcal{C}\!\mathit{oh}}
\def\Spacesstack{\mathcal{S}\!\mathit{paces}}
\def\Quotfunctor{\mathrm{Quot}}
\def\Hilbfunctor{\mathrm{Hilb}}
\def\Curvesstack{\mathcal{C}\!\mathit{urves}}
\def\Polarizedstack{\mathcal{P}\!\mathit{olarized}}
\def\Complexesstack{\mathcal{C}\!\mathit{omplexes}}
% \Pic is the operator that assigns to X its picard group, usage \Pic(X)
% \Picardstack_{X/B} denotes the Picard stack of X over B
% \Picardfunctor_{X/B} denotes the Picard functor of X over B
\def\Pic{\mathop{\mathrm{Pic}}\nolimits}
\def\Picardstack{\mathcal{P}\!\mathit{ic}}
\def\Picardfunctor{\mathrm{Pic}}
\def\Deformationcategory{\mathcal{D}\!\mathit{ef}}
\usepackage{cleveref}
\crefname{originalequation}{original equation}{original equations}
\Crefname{originalequation}{Original equation}{Original equations}
\makeatletter
\newlabel{equation-C-prime}{{0AL4}{source}{}{sourcepointer-1}{}}
\newlabel{equation-C-prime@cref}{{[originalequation][0][]0AL4}{[1][0][]source}{}{}{}}
\makeatother
\crefname{originalreference}{original reference}{original references}
\Crefname{originalreference}{Original reference}{Original references}
\makeatletter
\newlabel{lemma-NL-base-change}{{0GAG}{source}{}{sourcepointer-2}{}}
\newlabel{lemma-NL-base-change@cref}{{[originalreference][0][]0GAG}{[1][0][]source}{}{}{}}
\makeatother
\crefname{originallemma}{original lemma}{original lemmas}
\Crefname{originallemma}{Original lemma}{Original lemmas}
\makeatletter
\newlabel{more-algebra-lemma-base-change-property-ext-1-annihilated}{{0G9H}{source}{}{sourcepointer-3}{}}
\newlabel{more-algebra-lemma-base-change-property-ext-1-annihilated@cref}{{[originallemma][0][]0G9H}{[1][0][]source}{}{}{}}
\makeatother
\crefname{originallemma}{original lemma}{original lemmas}
\Crefname{originallemma}{Original lemma}{Original lemmas}
\makeatletter
\newlabel{more-algebra-lemma-surjection-property-ext-1-annihilated}{{0G9J}{source}{}{sourcepointer-4}{}}
\newlabel{more-algebra-lemma-surjection-property-ext-1-annihilated@cref}{{[originallemma][0][]0G9J}{[1][0][]source}{}{}{}}
\makeatother
\crefname{originallemma}{original lemma}{original lemmas}
\Crefname{originallemma}{Original lemma}{Original lemmas}
\makeatletter
\newlabel{more-algebra-lemma-two-term-base-change}{{0G9G}{source}{}{sourcepointer-5}{}}
\newlabel{more-algebra-lemma-two-term-base-change@cref}{{[originallemma][0][]0G9G}{[1][0][]source}{}{}{}}
\makeatother
\crefname{originalreference}{original reference}{original references}
\Crefname{originalreference}{Original reference}{Original references}
\makeatletter
\newlabel{remark-base-change}{{06FQ}{source}{}{sourcepointer-6}{}}
\newlabel{remark-base-change@cref}{{[originalreference][0][]06FQ}{[1][0][]source}{}{}{}}
\makeatother
\title{Annihilation of Ext\textasciicircum{}1 persists under base change}
\author{The Stacks Project Authors}
\date{stacks-project:restricted.tex:local-snapshot-20261012; GNU FDL 1.2 or later}
\begin{document}
\maketitle
\noindent\textit{Editorial note.} All delivered mathematical text is the author's own, copied from the pinned Stacks Project chapter file restricted.tex; no statement or proof step was rewritten, reordered or removed.\par\medskip
\noindent\textit{Editorial note.} The counted proof is the author's complete argument: the counted statement is the author lemma environment labelled lemma-zero-ext-1-after-modding-out; the proof environment immediately following it is its own proof\par\medskip
\noindent\textit{Editorial note.} The label lemma-zero-ext-1-after-modding-out names the counted statement, and the proof environment immediately following it in the source is delivered verbatim as its proof.\par\medskip
\noindent\textit{Editorial note.} The counted statement is carried by the author's own lemma environment in the chapter source.\par\medskip
\noindent\textit{Allowed standard prerequisites (editorial).}\begin{itemize}
\item the base change setup for pairs (A1, I1) to (A2, I2) of Noetherian adic rings (remark-base-change) and the base change B2 of a completion algebra B1 (equation-C-prime)
\item the base change map for naive cotangent complexes (lemma-NL-base-change): a map NL\_\{B1/A1\} tensor B2 B1 to NL\_\{B2/A2\} inducing an isomorphism on H\textasciicircum{}0 and a surjection on H\textasciicircum{}\{-1\}
\item the More on Algebra lemmas transferring vanishing/annihilation of Ext\textasciicircum{}1 under base change of a two-term complex and under a surjection on H\textasciicircum{}\{-1\} (more-algebra-lemma-two-term-base-change, more-algebra-lemma-base-change-property-ext-1-annihilated, more-algebra-lemma-surjection-property-ext-1-annihilated)
\end{itemize}\medskip
\ifdefined\hypertarget\hypertarget{sourcepointer-1}{}\fi
\noindent\textit{Source reference.} Original equation 0AL4: sources/stacks-restricted/pinned-original-source.tex, source lines 326--326. Prerequisite pointer invoked at proof line(s) 1019; it names a standard fact of the same chapter that the counted argument uses, and every step of the counted proof built on it is retained verbatim in the delivered body.\par\medskip
\ifdefined\hypertarget\hypertarget{sourcepointer-2}{}\fi
\noindent\textit{Source reference.} Original reference 0GAG: sources/stacks-restricted/pinned-original-source.tex, source lines 669--669. Prerequisite pointer invoked at proof line(s) 1029; it names a standard fact of the same chapter that the counted argument uses, and every step of the counted proof built on it is retained verbatim in the delivered body.\par\medskip
\ifdefined\hypertarget\hypertarget{sourcepointer-3}{}\fi
\noindent\textit{Source reference.} Original lemma 0G9H: sources/stacks-chow/tags.txt, source lines 20302--20302. Prerequisite pointer invoked at proof line(s) 1036; it names a standard fact of a sibling chapter that the counted argument uses, and every step of the counted proof built on it is retained verbatim in the delivered body.\par\medskip
\ifdefined\hypertarget\hypertarget{sourcepointer-4}{}\fi
\noindent\textit{Source reference.} Original lemma 0G9J: sources/stacks-chow/tags.txt, source lines 20304--20304. Prerequisite pointer invoked at proof line(s) 1037; it names a standard fact of a sibling chapter that the counted argument uses, and every step of the counted proof built on it is retained verbatim in the delivered body.\par\medskip
\ifdefined\hypertarget\hypertarget{sourcepointer-5}{}\fi
\noindent\textit{Source reference.} Original lemma 0G9G: sources/stacks-chow/tags.txt, source lines 20301--20301. Prerequisite pointer invoked at proof line(s) 1035; it names a standard fact of a sibling chapter that the counted argument uses, and every step of the counted proof built on it is retained verbatim in the delivered body.\par\medskip
\ifdefined\hypertarget\hypertarget{sourcepointer-6}{}\fi
\noindent\textit{Source reference.} Original reference 06FQ: sources/stacks-restricted/pinned-original-source.tex, source lines 426--426. Prerequisite pointer invoked at proof line(s) 1018; it names a standard fact of the same chapter that the counted argument uses, and every step of the counted proof built on it is retained verbatim in the delivered body.\par\medskip
\begin{lemma}
\label{lemma-zero-ext-1-after-modding-out}
Let $(A_1, I_1) \to (A_2, I_2)$ be as in
Remark \ref{remark-base-change} with $A_1$ and $A_2$ Noetherian.
Let $B_1$ be in (\ref{equation-C-prime}) for $(A_1, I_1)$.
Let $B_2$ be the base change of $B_1$. Let $f_1 \in B_1$
with image $f_2 \in B_2$.
If $\Ext^1_{B_1}(\NL_{B_1/A_1}^\wedge, N_1)$ is annihilated
by $f_1$ for every $B_1$-module $N_1$, then
$\Ext^1_{B_2}(\NL_{B_2/A_2}^\wedge, N_2)$ is annihilated
by $f_2$ for every $B_2$-module $N_2$.
\end{lemma}

\begin{proof}
By Lemma \ref{lemma-NL-base-change} there is a map
$$
\NL_{B_1/A_1} \otimes_{B_2} B_1 \to \NL_{B_2/A_2}
$$
which induces and isomorphism on $H^0$ and a surjection on $H^{-1}$.
Thus the result by More on Algebra, Lemmas
\ref{more-algebra-lemma-two-term-base-change},
\ref{more-algebra-lemma-base-change-property-ext-1-annihilated}, and
\ref{more-algebra-lemma-surjection-property-ext-1-annihilated}
the last two applied with the principal ideals $(f_1) \subset B_1$ and
$(f_2) \subset B_2$.
\end{proof}

\end{document}
